%
% moden.tex -- Die vier tiefsten Schwingungsmoden der kreisförmigen Membran
%
% (c) 2026 Gian Cavegn, OST Ostschweizer Fachhochschule
%
\documentclass[tikz]{standalone}
\usepackage{amsmath}
\usepackage{times}
\usepackage{txfonts}
\usetikzlibrary{arrows,intersections,math,calc}
\begin{document}
\def\skala{1}
\begin{tikzpicture}[>=latex,scale=\skala]
% Mode (0,1)
\begin{scope}
	\draw[thick] (0,0) circle (1.2);
	\node at (0,0) {$+$};
	\node at (0,-1.7) {$(0,1)$};
\end{scope}
% Mode (1,1)
\begin{scope}[xshift=3.1cm]
	\draw[thick] (0,0) circle (1.2);
	\draw[dashed] (0,-1.2) -- (0,1.2);
	\node at (-0.55,0) {$+$};
	\node at (0.55,0) {$-$};
	\node at (0,-1.7) {$(1,1)$};
\end{scope}
% Mode (2,1)
\begin{scope}[xshift=6.2cm]
	\draw[thick] (0,0) circle (1.2);
	\draw[dashed] (-0.849,-0.849) -- (0.849,0.849);
	\draw[dashed] (-0.849,0.849) -- (0.849,-0.849);
	\node at (0,0.62) {$+$};
	\node at (0,-0.62) {$+$};
	\node at (-0.62,0) {$-$};
	\node at (0.62,0) {$-$};
	\node at (0,-1.7) {$(2,1)$};
\end{scope}
% Mode (0,2)
\begin{scope}[xshift=9.3cm]
	\draw[thick] (0,0) circle (1.2);
	\draw[dashed] (0,0) circle (0.523);
	\node at (0,0) {$+$};
	\node at (0,0.86) {$-$};
	\node at (0,-1.7) {$(0,2)$};
\end{scope}
\end{tikzpicture}
\end{document}
